\subsection{真偏映射}

在本节中，X 与 Y 是局部紧拓扑空间。我们用 \(X'\)（相应地，\(Y'\)）表示在 X（相应地，Y）上添加一个无穷远点 \(\omega_{X}\)（相应地，\(\omega_{Y}\)）所得到的紧空间 (TG, I, p.~67--68)。我们把 \(X'\) 与 \(Y'\) 分别等同于 \(\mathsf{X}'(\mathcal{C}_{0}(\mathsf{X}))\) 与 \(\mathsf{X}'(\mathcal{C}_{0}(\mathsf{Y}))\) (I, p.~32, 推论 2)。

定义 1. --- 从 X 到 Y 的真偏映射是指 X 与 Y 之间的一个对应 \(f = (\Gamma, \mathrm{X}, \mathrm{Y})\) (E, II, p.~10, 定义 2)，它满足：

\begin{enumerate}
\def\labelenumi{(\roman{enumi})}
\item
  图像 \(\Gamma\) 是函数性的；
\item
  f 的定义域是 X 的一个开集 U；
\item
  从 U 到 Y 的映射 \(x\mapsto f(x)\) 是真映射。
\end{enumerate}

X 上的恒等映射是从 X 到 X 的真偏映射。设 Z 是一个局部紧拓扑空间，并设 \(f\)（相应地，\(g\)）是从 X 到 Y（相应地，从 Y 到 Z）的真偏映射。则复合对应 \(g \circ f\) (E, II, p.~11, 定义 6) 是从 X 到 Z 的真偏映射 (TG, I, p.~72, 命题 3 及 p.~73, 命题 5)。

引理 1. --- 对每个以 U 为定义域的、从 X 到 Y 的真偏映射 f，用 \(\widetilde{f}\) 表示由 \(\widetilde{f}(x) = f(x)\)（若 \(x \in \mathrm{U}\)）与 \(\widetilde{f}(x) = \omega_{\mathrm{Y}}\)（若 \(x \notin \mathrm{U}\)）定义的从 X' 到 Y' 的映射；它是连续的。

映射 \(f \mapsto \widetilde{f}\) 是从 X 到 Y 的真偏映射 \(f\) 所成的集合与从 X' 到 Y' 的连续映射 \(g\)（满足 \(g(\omega_{\mathrm{X}}) = \omega_{\mathrm{Y}}\)）所成的集合之间的一个双射。

设 \(f\) 是从 X 到 Y 的一个真偏映射，并设 U 是它的定义域。我们来证明映射 \(\widetilde{f}\) 是连续的。它在 U 的每一点都连续，因为 U 在 \(\mathrm{X}'\) 中是开的。我们来证明它在每一点 \(x\) ∈ \(\mathrm{X}' - \mathrm{U}\) 处也连续；这时有 \(\widetilde{f}(x) = \omega_{\mathrm{Y}}\)。设 V 是 \(\omega_{\mathrm{Y}}\) 在 \(\mathrm{Y}'\) 中的一个开邻域；我们来证明 \(\widetilde{f}^{-1}(\mathrm{V})\) 是 \(x\) 的一个邻域。根据拓扑空间 \(\mathrm{Y}'\) 的定义，可以假定 V 形如 \(Y^{\prime}-K\)，其中 K 是 Y 的一个紧子集。由于 f 定义了从 U 到 Y 的一个真映射，集合 \(f^{-1}(K)\) 在 U 中是紧的 (TG, I, p.~77, 命题 6)，因而在 \(X^{\prime}\) 中是紧的。它特别是 \(X^{\prime}\) 的一个闭子集，于是 \(\tilde{f}^{-1}(V)=X^{\prime}-f^{-1}(K)\) 是 \(X^{\prime}\) 的一个开子集，从而是 x 的一个邻域。

反之，设 \(g \colon X' \to Y'\) 是满足 \(g(\omega_{\mathrm{X}}) = \omega_{\mathrm{Y}}\) 的连续映射，并设 \(\Gamma_g \subset X' \times Y'\) 是它的图像。集合 \(U = X - \overline{g}^{-1}(\omega_{\mathrm{Y}})\) 在 \(X\) 中是开的。对应 \(f = (\Gamma_g \cap (U \times Y), X, Y)\) 是从 \(X\) 到 \(Y\) 的一个真偏映射，它满足 \(\widetilde{f} = g\)，并且是唯一的一个。

我们将把从 X 到 Y 的真偏映射与从 \(X'\) 到 \(Y'\) 中的、把 \(\omega_{X}\) 映到 \(\omega_{Y}\) 的连续映射视为同一。特别地，从 X 到 Y 的真映射就是以 X 为定义域的真偏映射；它们与从 \(X'\) 到 \(Y'\) 中的、满足 \(f^{-1}(\omega_{\mathrm{Y}}) = \{\omega_{\mathrm{X}}\}\) 的连续映射 f 视为同一。若 X 是紧的，这些映射无非就是从 X 到 Y 的连续映射。

设 A 是一个复交换 Banach 代数。回顾 \(X'(A)\) 与在 \(X(A)\) 上添加一个无穷远点所得到的紧空间视为同一 (I, p.~29, 推论)。

命题 2. --- 设 A 与 B 是复交换 Banach 代数。对每个代数同态 \(\pi\colon\) A \(\to\) B，映射 \(\mathsf{X}'(\pi)\) 是从 \(\mathsf{X}(\mathsf{B})\) 到 \(\mathsf{X}(\mathsf{A})\) 的真偏映射。

事实上，\(\mathsf{X}'(\pi)\) 是从 \(\mathsf{X}'(\mathrm{B})\) 到 \(\mathsf{X}'(\mathrm{A})\) 的连续映射 (I, p.~10)。\(\mathsf{X}'(\mathrm{B})\)（相应地，\(\mathsf{X}'(\mathrm{A})\)）的无穷远点是零特征标，并且我们有 \(\mathsf{X}'(\pi)(0)=0\)。

命题 3. --- a) 对每个从 X 到 Y 的真偏映射 \(\varphi\)，映射 \(f \mapsto f \circ \varphi\)（从 \(\mathcal{C}(\mathrm{Y}')\) 到 \(\mathcal{C}(\mathrm{X}')\)）诱导出一个代数态射 \(\varphi^*\)（从 \(\mathcal{C}_0(\mathrm{Y})\) 到 \(\mathcal{C}_0(\mathrm{X})\)）；

\begin{enumerate}
\def\labelenumi{\alph{enumi})}
\setcounter{enumi}{1}
\tightlist
\item
  映射 \(\varphi \mapsto \varphi^{*}\) 是从 X 到 Y 的真偏映射所成的集合到由 \(\mathcal{C}_0(\mathrm{Y})\) 到 \(\mathcal{C}_0(\mathrm{X})\) 的代数态射所成的集合上的一个双射。其逆双射是映射 \(\pi \mapsto \mathsf{X}'(\pi)\)。
\end{enumerate}

我们来证明 a)。设 \(\varphi\) 是从 X 到 Y 的真偏映射，把它与满足 \(\varphi(\omega_{\mathrm{X}}) = \omega_{\mathrm{Y}}\) 的、从 X' 到 Y' 的连续映射视为同一。对 \(f \in \mathcal{C}_0(\mathrm{Y})\)，有 \((f \circ \varphi)(\omega_{\mathrm{X}}) = f(\omega_{\mathrm{Y}}) = 0\)，因而映射 \(\varphi^*\) 是良定义的。它是一个代数态射。

我们来证明 \(\mathsf{X}'(\varphi^{*})\) 与 \(\varphi\) 视为同一。设 \(x \in X\)。对每个函数 \(f \in \mathcal{C}_{0}(\mathrm{Y})\)，特征标 \(\mathsf{X}'(\varphi^{*})(\mathrm{ev}_{x})\) 把复数

\[
\left(\operatorname{ev} _ {x} \circ \varphi^ {*}\right) (f) = \operatorname{ev} _ {x} (f \circ \varphi) = f (\varphi (x)),
\]

与 f 相联系，因而 \(\mathsf{X}'(\varphi^{*})(\mathrm{ev}_{x}) = \mathrm{ev}_{\varphi(x)}\)。这就证明了该断言。

反之，设 \(\pi\colon\mathcal{C}_{0}(\mathrm{Y})\to\mathcal{C}_{0}(\mathrm{X})\) 是一个代数态射。我们来证明 \(\mathsf{X}^{\prime}(\pi)^{*}=\pi\)。设 \(f\in\mathcal{C}_{0}(\mathrm{Y})\)，并记 \(g=\mathsf{X}^{\prime}(\pi)^{*}(f)\in\mathcal{C}_{0}(\mathrm{X})\)。对每个 \(x\in\mathsf{X}\)，有

\[
g (x) = \left(f \circ \mathrm{X} ^ {\prime} (\pi)\right) (x) = \operatorname{ev} _ {\mathrm{X} ^ {\prime} (\pi) (x)} (f) = \left(\operatorname{ev} _ {x} \circ \pi\right) (f) = \pi (f) (x),
\]

因为 \(\mathsf{X}^{\prime}(\pi)\) 满足 \(\mathrm{ev}_{\mathsf{X}^{\prime}(\pi)(x)} = \mathrm{ev}_{x} \circ \pi\)。于是我们有 \(g = \pi(f)\)，由此得出 \(\mathsf{X}^{\prime}(\pi)^{*} = \pi\)。

以完全类似的方式，我们有：

命题 4. --- 假定 X 与 Y 是紧的。把空间 X（相应地，空间 Y）与 X( \(\mathcal{C}(X)\) )（相应地，X( \(\mathcal{C}(Y)\) )）视为同一 (I, p.~33, 推论 3)。

\begin{enumerate}
\def\labelenumi{\alph{enumi})}
\item
  对每个连续映射 \(\varphi\colon\mathrm{X}\to\mathrm{Y}\)，映射 \(\varphi^{*}\colon f\mapsto f\circ\varphi\) 是从 \(\mathcal{C}(\mathrm{Y})\) 到 \(\mathcal{C}(\mathrm{X})\) 的代数态射；
\item
  映射 \(\varphi \mapsto \varphi^{*}\) 与 \(\pi \mapsto X(\pi)\) 是从 \(X\) 到 \(Y\) 的连续映射所成的集合与从 \(\mathcal{C}(Y)\) 到 \(\mathcal{C}(X)\) 的代数态射所成的集合之间互逆的双射。
\end{enumerate}

注. --- *用范畴论的语言，前面的结果可解释如下。设 G 是这样一个范畴：其对象为局部紧拓扑空间，其态射为真偏映射。函子 X ↦ \(\mathcal{C}_0(X)\) 是从范畴 G 到复交换 Banach 代数范畴的一个反变的全忠实函子。此外，A ↦ X(A) 是从复交换 Banach 代数范畴到范畴 G 的一个反变函子。若对每个局部紧拓扑空间 X 联系同胚

\[
\operatorname{ev} \colon X \to X (\mathscr {C} _ {0} (X)),
\]

则得到从范畴 G 的恒等函子到复合函子 \(X \mapsto X(\mathcal{C}_{0}(X))\) 的一个同构。

复合函子 \(A \mapsto \mathcal{C}_{0}(X(A))\) 并不同构于复交换 Banach 代数范畴的恒等函子（参见 I, p.~36, 例 2 及 I, p.~155, 习题 2）。不过，对于交换星代数（C*-代数），我们将看到一个这种类型的命题 (I, p.~107, 第 5 小节）。*

